The aim for this section is to introduce some statistical notions that will be used in order to solve the exercises in this homework.
We begin exploring the concept of a \textit{statistical pivot}. A statistical pivot (or \emph{pivotal quantity}) is a function of both the sample data and one or more unknown parameters whose \textbf{probability distribution does not depend on those parameters}. 
\begin{definition}[Statistical Pivot]

Formally, if $X = (X_1, \dots, X_n)$ is a random sample from a distribution depending on a parameter $\theta$, a statistic $T(X, \theta)$ is called a \emph{pivot} if the distribution of $T$ is the same for all possible values of $\theta$.
\end{definition}

In simpler terms, a pivot is a clever combination of our data $X$ and the unknown parameters $\theta$ that “cancels out” the dependence on those parameters, leaving something whose behavior is fully known. As we'll explore throughout the worked exercises,
pivotal quantities are extremely useful in statistical inference because they allow us to construct \textbf{confidence intervals} and perform \textbf{hypothesis tests} without knowing the true value of the underlying parameter.

\begin{remark} 
By using a pivot, we can make probability statements that remain valid for any possible value of $\theta$.
\end{remark}

\vspace{0.2cm}

We now explore the concept of \textit{equivariant estimators}. Equivariant estimators are statistical estimators that maintain a consistent relationship with the data under certain transformations.

\begin{definition}[Equivariant estimators]
    A parameter estimator $g(y)$ is said to be \textbf{equivariant} if it satisfies:
    $$g(a + by) = g(a + by_1, \ldots, a + by_1) = a + g(by_1, \ldots, by_n) = a + bg(y)$$
    for all $a \in k$ and $b > 0$.
\end{definition}

It is useful to recall the following remark regarding the equivariance properties of estimators in a location-scale model.

\begin{remark}
For estimators in a location–scale model 
\[
Y_i = \eta + \tau V_i, \quad i=1,\dots,n,
\]
it is natural to require that they behave consistently under changes of location and scale. 
A \textbf{location estimator} $m(\cdot)$ is said to be \emph{location–scale equivariant} if
\[
m(\eta + \tau v_1, \ldots, \eta + \tau v_n) = \eta + \tau\, m(v_1, \ldots, v_n),
\]
that is, translating and rescaling the data simply translates and rescales the estimator in the same way. Similarly, a \textbf{scale estimator} $s(\cdot)$ is said to be \emph{location invariant} and \emph{scale equivariant} if
\[
s(\eta + \tau v_1, \ldots, \eta + \tau v_n) = |\tau|\, s(v_1, \ldots, v_n).
\]
These properties ensure that both estimators respond to shifts and rescalings of the data, preserving the structure of the underlying model.
\end{remark}

\vspace{0.2cm}

We now introduce the notion of quantiles, which will be useful in the exercises.
\begin{definition}[Quantiles]
    For a random variable $X$ with cumulative distribution function (CDF) $F_X(x)$, the quantile $q_p$ is the value such that:
    $$\mathbb{P}(X \leq q_p) = p$$
    i.g. $q_{0.5}$ is the median of $X$. This is just a fancy way of saying `the point where the CDF hits a certain probability'.    
\end{definition}

\vspace{0.2cm}

We briefly recall the definition of median and interquartile range (IQR), which will be used in the exercises.
\begin{definition}[Median]
The median of a dataset is the \textbf{middle value} when the data points are arranged in ascending order. Thus, it may be thought of as a quantile $q_{0.5}$ such
that
$$\mathbb{P}(q_i \leq q_{0.5}) = 0.5$$
\begin{itemize}
    \item If $n$ is odd, the median is the value at position $\frac{n+1}{2}$.
    \item If $n$ is even, the median is the average of the values at positions $\frac{n}{2}$ and $\frac{n}{2} + 1$.
\end{itemize} 
\end{definition}

Now, we define the interquartile range (IQR) as follows.

\begin{definition}[Interquartile Range]
    The interquartile range (IQR) is a measure of statistical dispersion, or how spread out the values in a dataset are. It is defined as the difference between the third quartile ($Q_3$) and the first quartile ($Q_1$):
    $$\text{IQR} = Q_3 - Q_1$$
    where:
    \begin{itemize}
        \item $Q_1$ (the first quartile) is the median of the lower half of the dataset (the 25th percentile i.e. the $q_{.25}$ quantile!).
        \item $Q_3$ (the third quartile) is the median of the upper half of the dataset (the 75th percentile or $q_{.75}$ quantile).
    \end{itemize}
\end{definition}

Note that percentiles and quantiles are very closely related concepts. Percentiles are specific types of quantiles that divide a dataset into 100 equal parts. Thus, 
percentiles can be viewed as quantiles expressed in terms of percentages.